\documentclass[10pt,aspectratio=169]{beamer}
\usetheme{default}

\usepackage{booktabs}
\usepackage{listings}
\usepackage{xcolor}

\definecolor{codebg}{HTML}{F5F5F5}
\definecolor{codecmt}{HTML}{6A737D}
\definecolor{codekw}{HTML}{0044AA}
\definecolor{codestr}{HTML}{B36000}

\lstdefinelanguage{Rsimple}{
  keywords={brm, gaussian, normal, student_t, s, t2, se, condition,
            participant, prior, bs, by, k, d, family, c, eeg_mean,
            eeg_sd, xproj, yproj, time},
  keywordstyle=\color{codekw}\bfseries,
  commentstyle=\color{codecmt}\itshape,
  stringstyle=\color{codestr},
  morecomment=[l]{\#},
  morestring=[b]",
  morestring=[b]',
  sensitive=true,
}

\lstset{
  language=Rsimple,
  basicstyle=\ttfamily\footnotesize,
  backgroundcolor=\color{codebg},
  frame=single,
  rulecolor=\color{codebg},
  framesep=4pt,
  showstringspaces=false,
  breaklines=true,
  columns=fullflexible,
}

\title{Spatio-temporal BGAM for EEG}
\subtitle{Pilot validation on simulated data}
\author{Jiahui Fan}
\date{2026-05-11}

\begin{document}

\maketitle

% ===================================================================
% Frame 1 - Setup
% ===================================================================
\begin{frame}{Setup}
  \begin{itemize}
    \item \textbf{Pilot data} (simulated): 5 participants
          $\times$ 20 trials/condition $\times$ 61 electrodes
          $\times$ 51 time points (generated at 500 Hz,
          downsampled $\times 5 \to 100$ Hz).
    \item \textbf{Two conditions}: \texttt{cond1} (noise only),
          \texttt{cond2} (noise + signal).
    \item \textbf{Injected signal}: spatial Gaussian centred at CZ
          ($\sigma_{\text{space}} = 0.5$), amplitude $1.5\,\mu\mathrm{V}$;
          per-participant onset
          $\tau_p \sim \mathrm{Uniform}[150, 170]\,\mathrm{ms}$,
          ERP duration fixed at $182\,\mathrm{ms}$
          $\Rightarrow$ offset $\in [332, 352]\,\mathrm{ms}$
          (central true window $160$--$342\,\mathrm{ms}$).
  \end{itemize}
\end{frame}

% ===================================================================
% Frame 1b - Data generation
% ===================================================================
\begin{frame}{Data generation}
  \textbf{Signal} (cond2 only):
  $y_c(t) = A \cdot w_c \cdot \mathrm{erp}(t - \tau_p)$
  \begin{itemize}\setlength{\itemsep}{0pt}
    \item Spatial weight: $w_c = \exp(-d_c^2 / 2\sigma_s^2)$,
          $\sigma_s = 0.5$, centred at CZ;
          $d_c$ = standardised Euclidean distance
          from electrode $c$ to CZ.
    \item ERP shape: Gaussian curve (width $= 182$ ms).
    \item Amplitude $A = 1.5\,\mu\mathrm{V}$;
          $\tau_p \sim \mathrm{Uniform}[150, 170]\,\mathrm{ms}$.
  \end{itemize}

  \vspace{0.3em}
  \textbf{Noise} (both conditions, independent draws):
  \begin{itemize}\setlength{\itemsep}{0pt}
    \item Temporal: $1/f$ pink noise
          (Yeung et al.\ 2004 EEG spectrum).
    \item Spatial: correlated $\mathcal{N}(0, \Sigma)$,
          $\Sigma_{ij} = \exp(-d_{ij}^2 / 2\rho^2)$, $\rho = 0.15$.
    \item Temporal var $=$ spatial var $= 0.5$.
  \end{itemize}
\end{frame}

% ===================================================================
% Frame 2 - Model
% ===================================================================
\begin{frame}[fragile]{Model: BGAM (v4\_k12)}
\begin{lstlisting}
brm(
  eeg_mean | se(eeg_sd) ~
    condition +
    (1 + condition | participant) +
    s(time, bs = "tp", k = 20, by = condition) +
    t2(xproj, yproj, time,
       by = condition,
       bs = c("tp", "tp"),
       k  = c(10, 12), d = c(2, 1)),
  family = gaussian(),
  prior  = c(normal(0, 1), student_t(3, 0, 1) x2))
\end{lstlisting}

  \vspace{0.3em}
  \begin{itemize}
    \item \texttt{| se(eeg\_sd)}: measurement-error model
          (per-cell trial SD).
    \item \texttt{t2(xproj, yproj, time)}:
          joint 3D smooth, $k = (10,12)$, $d = (2,1)$.
  \end{itemize}
\end{frame}

% ===================================================================
% Frame 3 - Fitting & convergence
% ===================================================================
\begin{frame}{Fitting \& convergence}
  \textbf{Sampler}: 4 chains $\times$ 3000 iter (warmup 1000),
  threading $= 4$/chain, \texttt{adapt\_delta} $= 0.99$,
  \texttt{max\_treedepth} $= 13$, \texttt{cmdstanr} backend.
  $\sim$14 h wall time on LIDO HPC (16 cores).

  \vspace{0.6em}
  \begin{center}
  \begin{tabular}{lrr}
    \toprule
    Metric                & Value           & Threshold \\
    \midrule
    Divergent transitions & \textbf{0}      & 0          \\
    Max treedepth hits    & \textbf{0}      & $< 5\%$    \\
    Max $\hat{R}$         & \textbf{1.0026} & $< 1.01$   \\
    Min Bulk\_ESS         & \textbf{2{,}435} & $> 400$   \\
    Min Tail\_ESS         & \textbf{4{,}113} & $> 400$   \\
    \bottomrule
  \end{tabular}
  \end{center}

  \vspace{0.5em}
  All diagnostics pass; chains mix cleanly.
\end{frame}

% ===================================================================
% Frame 4 - Recovery (numbers)
% ===================================================================
\begin{frame}{Recovery (single fit)}
  \textbf{Temporal recovery} (seed electrode CZ):
  \begin{center}
  \begin{tabular}{lccc}
    \toprule
                  & True   & Estimated & Bias \\
    \midrule
    Onset (ms)    & 160    & 150       & $-10$ \\
    Offset (ms)   & 342    & 350       & $+8$  \\
    \bottomrule
  \end{tabular}
  \end{center}

  \vspace{0.6em}
  \textbf{Spatial recovery} (peak time 250 ms):
  \begin{itemize}
    \item Top channel = \textbf{CZ} (= seed) \checkmark.
    \item CZ posterior \texttt{cond2 - cond1}:
          $+0.554\,\mu\mathrm{V}$ \; (95\% CrI [0.459, 0.640]).
    \item Top 10 = CZ + its 9 immediate 10--20 neighbours.
  \end{itemize}

  \vspace{0.4em}
  \textbf{Empirical vs.\ model, two views}:
  \begin{center}
  \footnotesize
  \begin{tabular}{lcc}
    \toprule
                & \textbf{Marginal} (global avg) & \textbf{Peak} (CZ @ 250 ms) \\
    \midrule
    Empirical   & $0.055\,\mu\mathrm{V}$ & $\sim 1.5\,\mu\mathrm{V}$ \\
    BGAM model  & $0.056\,\mu\mathrm{V}$ & $0.55\,\mu\mathrm{V}$     \\
    \midrule
    Agreement   & \checkmark $\;$global mean unbiased
                & $\sim$37\%: peak amplitude shrunk \\
    \bottomrule
  \end{tabular}
  \end{center}
\end{frame}

% ===================================================================
% Frame 4 - Figure
% ===================================================================
\begin{frame}{Spatio-temporal recovery (peak time 250 ms)}
  \begin{center}
    \includegraphics[height=0.78\textheight]{../results/figures/topoplot_4rows_gt_emp_vs_model.pdf}
  \end{center}

  \vspace{-0.3em}
  \footnotesize
  Row 2/3 look pale because BGAM shrinks the peak to
  $\sim$37\% of the true amplitude ($0.55$ vs.\ $1.5\,\mu\mathrm{V}$),
  a known smoothing-penalty effect.
\end{frame}

% ===================================================================
% Frame 5 - Method comparison
% ===================================================================
\begin{frame}{Method comparison (30-sim Monte Carlo)}
  Same 3-step rule (mask $\to$ blob filter $\to$ largest blob).
  Only Step 1 differs: raw $p$ (univariate $t$-test, $\alpha = 0.1$),
  cluster perm.\ (MNE 2D \texttt{spatio\_temporal\_cluster\_1samp\_test}),
  BGAM\_2D (posterior odds $> 10$).

  \vspace{0.3em}
  \begin{center}
  \footnotesize
  \begin{tabular}{lrrrr}
    \toprule
    Method            & Bias & MAE & RMSE & Variance \\
    \midrule
    \multicolumn{5}{l}{\textit{Onset (ms)}} \\
    raw $p$           & $+9.67$  & \textbf{10.33} & 12.25          & 58.51          \\
    \textbf{BGAM\_2D} & $-10.67$ & 10.67          & \textbf{11.55} & \textbf{20.23} \\
    cluster perm.     & $+15.67$ & 15.67          & 17.22          & 52.99          \\
    \addlinespace
    \multicolumn{5}{l}{\textit{Offset (ms)}} \\
    \textbf{BGAM\_2D} & $+12.33$ & \textbf{12.47} & \textbf{13.78} & \textbf{39.20}  \\
    raw $p$           & $-6.00$  & 12.53          & 15.28          & 204.14          \\
    cluster perm.     & $-18.00$ & 18.53          & 20.03          & 80.00           \\
    \bottomrule
  \end{tabular}
  \end{center}

  \vspace{0.3em}
  \footnotesize
  \textbf{Joint spatio-temporal $F_1$} (mean $\pm$ sd over 30 sims):
  BGAM\_2D $= 0.799 \pm 0.007$,
  cluster perm.\ $= 0.514 \pm 0.020$,
  raw $p = 0.462 \pm 0.014$.
\end{frame}

% ===================================================================
% Backup slide - Joint ST F1 definition
% ===================================================================
\appendix

\begin{frame}{Backup: Joint spatio-temporal $F_1$}
  \textbf{Standard $F_1$ on the $(\text{time} \times \text{channel})$ grid}:
  \[
    F_1 = \frac{2 \cdot \text{precision} \cdot \text{recall}}
                {\text{precision} + \text{recall}},
    \quad
    \text{precision} = \frac{TP}{TP + FP},
    \quad
    \text{recall} = \frac{TP}{TP + FN}.
  \]

  \vspace{0.4em}
  \textbf{Truth mask} (from ground truth):
  \begin{itemize}\setlength{\itemsep}{0pt}
    \item Active electrodes: \texttt{spatial\_weight}$_c > 0.1$
          ($\sim$25 channels around CZ).
    \item Active times: $160 \leq t \leq 342$ ms (19 time points,
          using the fixed central onset/offset).
    \item $475$ TRUE cells, $15.3\%$ of the grid.
  \end{itemize}

  \vspace{0.4em}
  \textbf{Prediction mask}: each method's blob-filtered significant mask
  (after Steps 1--2 of the 3-step rule).

  \vspace{0.4em}
  Computed per simulation; reported as mean $\pm$ sd across 30 MC sims.
\end{frame}

% ===================================================================
% Backup slide - why ALPHA = 0.1 (sign-flip permutation derivation)
% ===================================================================
\begin{frame}{Backup: why $\alpha = 0.1$ at $n = 5$?}
  \textbf{Sign-flip permutation test}: flip the sign of each participant's
  (cond2 $-$ cond1) value, or keep it.

  \vspace{0.3em}
  \begin{itemize}\setlength{\itemsep}{0pt}
    \item Total flip configurations: $2^n$
          (each of $n$ participants: flip or not).
    \item $n = 5 \Rightarrow 2^5 = 32$ configurations.
    \item \textbf{Two-sided test}: only $|\text{test statistic}|$
          matters, so the two most extreme configurations
          (most positive and most negative) both count.
    \item \textbf{Smallest possible $p$-value}:
          $\displaystyle p_{\min} = \frac{2}{2^n} = \frac{1}{2^{n-1}}.$
    \item $n = 5 \Rightarrow p_{\min} = \tfrac{1}{16} = 0.0625$.
  \end{itemize}

  \vspace{0.3em}
  $\alpha = 0.05$ is mathematically unreachable
  $(0.0625 > 0.05)$; we need $0.0625 < 0.1$.

  \vspace{0.3em}
  Once $n \geq 10$: $p_{\min} = \tfrac{1}{2^9} \approx 0.002$,
  so $\alpha = 0.05$ becomes attainable.
\end{frame}

\end{document}
